% ===========================================================
%  第 2 章 余子式和代数余子式的计算 —— 高等代数【深化】拓展框（主控撰写）
%  对应 OUTLINE.md §3「深化清单」2.1 – 2.3
%  用法：\begin{fdeep}{标题} 灰底框；由主控插入正文相应考点旁
%  注意：本文件第 1 行以下的块是"手写块"，提取脚本只认行首的 \begin{fdeep}
% ===========================================================

% ---------- 2.1 伴随矩阵的完整理论（对应 OUTLINE 2.1）----------
\begin{fdeep}{伴随矩阵：一个恒等式定全局}
\textbf{定义}：设 $A=(a_{ij})_{n\times n}$，$A_{ij}$ 为 $a_{ij}$ 的代数余子式，则
\fml{A^{*}=\begin{pmatrix}A_{11}&A_{21}&\cdots&A_{n1}\\ A_{12}&A_{22}&\cdots&A_{n2}\\ \vdots&\vdots&&\vdots\\ A_{1n}&A_{2n}&\cdots&A_{nn}\end{pmatrix}}
\textbf{注意转置}：$A^{*}$ 的第 $i$ \textbf{行}是 $A$ 第 $i$ \textbf{列}元素的代数余子式。这一条错位是最常见的失分点。

\medskip
\textbf{唯一需要记住的恒等式}：
\fml{AA^{*}=A^{*}A=|A|E}
\textbf{它从哪来}：$(AA^{*})_{ij}=\sum_{k=1}^{n}a_{ik}A_{jk}$，即"第 $i$ 行元素乘第 $j$ 行元素的代数余子式"。
由行列式展开定理，
\fml{\sum_{k=1}^{n}a_{ik}A_{jk}=\begin{cases}|A|,& i=j\\ 0,& i\ne j\end{cases}}
——$i=j$ 时是"按第 $i$ 行展开"，$i\ne j$ 时相当于把第 $i$ 行换成第 $j$ 行（两行相同，行列式为 $0$）。
所以 $AA^{*}=|A|E$ 不是新结论，而是\textbf{展开定理的直接改写}。

\medskip
\textbf{由它一网打尽全部常用结论}（$n\ge2$ 时）：
\begin{itemize}[leftmargin=2.2em,itemsep=0.22em]
\item $|A|\ne0$ 时，$A^{-1}=\dfrac{1}{|A|}A^{*}$ —— 这就是\textbf{用伴随矩阵求逆}的依据；
\item $|A^{*}|=|A|^{\,n-1}$ —— 对 $AA^{*}=|A|E$ 两边取行列式，$|A|\,|A^{*}|=|A|^{n}$，
      再对 $|A|=0$ 的情形作连续性/按定义单独论证即得；
\item $(A^{*})^{*}=|A|^{\,n-2}A$（$n\ge3$）；
\item $(kA)^{*}=k^{\,n-1}A^{*}$，$(A^{\mathrm{T}})^{*}=(A^{*})^{\mathrm{T}}$；
\item $A^{*}=|A|A^{-1}$ 说明\textbf{$A^{*}$ 与 $A$ 是"同类"的}：$A$ 对称则 $A^{*}$ 对称，$A$ 可逆则 $A^{*}$ 可逆。
\end{itemize}
【反哺点】考纲"矩阵"一节要求"理解伴随矩阵的概念，会用伴随矩阵求逆矩阵"。
把这一串结论全部挂在 $AA^{*}=|A|E$ 这一个等式上，就不必逐条背——
考场上先写出 $AA^{*}=|A|E$，再问"我要的是哪一项"，绝大多数伴随矩阵题都能现场推出来。
\end{fdeep}

% ---------- 2.1b 秩公式 r(A*) 的三分（对应 OUTLINE 2.1）----------
\begin{fdeep}{$r(A^{*})$ 的三分判别：$n$ / $1$ / $0$}
设 $A$ 为 $n$ 阶矩阵（$n\ge2$），则
\fml{r(A^{*})=\begin{cases}n,& r(A)=n\\ 1,& r(A)=n-1\\ 0,& r(A)\le n-2\end{cases}}
\textbf{三段证明思路}：
\begin{enumerate}[leftmargin=2.2em,itemsep=0.25em]
\item $r(A)=n$（即 $|A|\ne0$）：由 $|A^{*}|=|A|^{\,n-1}\ne0$，故 $A^{*}$ 可逆，$r(A^{*})=n$；
\item $r(A)=n-1$：此时 $A$ 存在一个 $n-1$ 阶非零子式，故\textbf{至少有一个 $A_{ij}\ne0$}，即 $A^{*}\ne O$，$r(A^{*})\ge1$；
      又由 $AA^{*}=|A|E=O$ 知 $A^{*}$ 的每一列都是 $Ax=0$ 的解向量。
      而 $r(A)=n-1$ 时 $Ax=0$ 的基础解系只含 $n-r(A)=1$ 个向量，
      故 $A^{*}$ 的各列都与同一个向量成比例，$r(A^{*})\le1$。合起来 $r(A^{*})=1$；
\item $r(A)\le n-2$：此时 $A$ 的所有 $n-1$ 阶子式全为 $0$，故全部 $A_{ij}=0$，即 $A^{*}=O$，$r(A^{*})=0$。
\end{enumerate}
\textbf{第 ② 段的思路最值得记}：它把"伴随矩阵的秩"转化为\textbf{"解空间的维数只有 1"}\allowbreak——
这正好复用第 5 章"基础解系含 $n-r$ 个向量"的结论。

【反哺点】这是考研选择题的高频送分点，标准问法是"已知 $r(A)=n-1$，求 $r(A^{*})$"。
另外它给出一个反向用法：\textbf{由 $r(A^{*})$ 反推 $r(A)$}——见到 $A^{*}=O$ 就意味着 $r(A)\le n-2$。
\end{fdeep}

% ---------- 2.2 代数余子式线性组合的构造技巧（对应 OUTLINE 2.2）----------
\begin{fdeep}{代数余子式的和怎么求：把"配系数"翻译成"构造行列式"}
本讲的核心手法是：\textbf{代数余子式前的系数，就是替换掉的那一行（列）的元素}。
\fml{\sum_{k=1}^{n}k_{k}A_{ik}=\begin{vmatrix}
a_{11}&\cdots&a_{1n}\\
\vdots&&\vdots\\
k_{1}&\cdots&k_{n}&\leftarrow\text{第 }i\text{ 行被换成 }k_{1},\dots,k_{n}\\
\vdots&&\vdots\\
a_{n1}&\cdots&a_{nn}
\end{vmatrix}}
也就是说：\textbf{$A_{i1},A_{i2},\dots,A_{in}$ 只由 $A$ 中\emph{去掉第 $i$ 行}后的部分决定}，
与第 $i$ 行元素本身无关。既然无关，就可以\textbf{自由地把第 $i$ 行换成任何想要的行}——
只要那个新行列式好算，就算出了这组代数余子式的线性组合。

\medskip
\textbf{由此得到一套"求余子式之和"的固定套路}：
\begin{enumerate}[leftmargin=2.2em,itemsep=0.25em]
\item 要把 $\sum A_{i1}$（第 $i$ 行代数余子式之和）求出来，
      就在 $\sum_k a_{ik}A_{ik}=|A|$ 里\textbf{把 $a_{ik}$ 全部换成 $1$}，得
      \fml{\sum_{k=1}^{n}A_{ik}=\begin{vmatrix}\cdots\\ 1&1&\cdots&1\\ \cdots\end{vmatrix}}
      即"把第 $i$ 行整行换成 $1$"之后的行列式；
\item 要求 $\sum_{i,j}A_{ij}$（\textbf{全部}代数余子式之和），就把它看成
      $\sum_i\left(\sum_j A_{ij}\right)$：对每个 $i$ 都把第 $i$ 行换成全 $1$，再全部相加。
      更快的一步到位写法是左乘全 $1$ 矩阵：
      \fml{\sum_{i=1}^{n}\sum_{j=1}^{n}A_{ij}=\mathbf{1}^{\mathrm{T}}A^{*}\mathbf{1}=\mathbf{1}^{\mathrm{T}}\bigl(|A|A^{-1}\bigr)\mathbf{1}}
      （$|A|\ne0$ 时），于是问题变成\textbf{先求 $A^{-1}$ 或直接算 $A^{*}$ 的一个二次型}；
\item 若系数不是全 $1$ 而是 $k_1,\dots,k_n$，则\textbf{逐位替换}即可，不必重新展开。
\end{enumerate}
【反哺点】9 讲把这一招按"用行列式 / 用矩阵 / 用特征值"三条路线讲，
但你只要记住一句话——\textbf{"代数余子式的系数 = 被换掉的那一行"}，
三条路线其实是同一件事的三种算法。考场上先写出 $\sum a_{ik}A_{ik}=|A|$，
再问"我要的系数是什么"，把系数填进被替换行，题就变成算一个行列式。
\end{fdeep}

% ---------- 2.3 |A*| 与伴随矩阵的"多项式关系"（对应 OUTLINE 2.3）----------
\begin{fdeep}{$|A^{*}|=|A|^{\,n-1}$ 为什么成立，以及它的两个推论}
\textbf{推导只需一行}：在 $AA^{*}=|A|E$ 两边取行列式，用 $|kE|=k^{n}$，得
\fml{|A|\cdot|A^{*}|=\bigl||A|E\bigr|=|A|^{\,n}\;\Longrightarrow\;|A^{*}|=|A|^{\,n-1}\qquad(|A|\ne0)}
\textbf{$|A|=0$ 的情形}要单独说：此时 $r(A)\le n-1$。
若 $r(A)\le n-2$ 则 $A^{*}=O$，$|A^{*}|=0=|A|^{n-1}$ 成立；
若 $r(A)=n-1$ 则 $r(A^{*})=1$，$A^{*}$ 的各行成比例，仍得 $|A^{*}|=0=|A|^{n-1}$。
（这里正好复用了「$r(A^{*})$ 三分」那一条，两条结论互为支撑。）

\medskip
\textbf{两个常被忽略的推论}：
\begin{itemize}[leftmargin=2.2em,itemsep=0.22em]
\item $(A^{*})^{*}=|A|^{\,n-2}A$（$n\ge3$）。
      思路：对 $A^{*}$ 用同一公式，$(A^{*})^{*}=|A^{*}|(A^{*})^{-1}=|A|^{\,n-1}\cdot\dfrac{A}{|A|}=|A|^{\,n-2}A$；
      \textbf{注意 $n=2$ 时此式退化为 $(A^{*})^{*}=A$}，不要在 $n=2$ 的题里套错指数。
\item $|A^{*}|=|A|^{\,n-1}$ 的\textbf{反向用法}：见到 $|A^{*}|$ 与 $|A|$ 同时出现，
      可以直接把次数配平。例如由 $|A^{*}|=8$ 且 $n=3$ 立即得 $|A|^{2}=8$，$|A|=\pm2\sqrt2$ 需再定号。
\end{itemize}
【反哺点】考纲要求"会用伴随矩阵求逆矩阵"，而 $|A^{*}|$、$(A^{*})^{*}$ 这类量
正是填空题的常见落点。它们没有独立记忆的必要——全部挂在 $AA^{*}=|A|E$ 上现推即可，
但\textbf{$n=2$ 与 $n\ge3$ 的指数差别}必须记牢。
\end{fdeep}
